\begin{textsample}{NEG Dim 7 – global\_south\_2024\_12 – Score -4.00 – global\_south\_2024\_12/118656078.txt}  \label{ex:f7_neg_300}
About 50 people picketed outside the Samson Centre in Cape Town on Monday to mark the UN International Day of \textbf{Commemoration} and Dignity of the Victims of the Crime of Genocide.
During the protest, two women wearing white were covered in red \textbf{paint} to symbolise genocides.
Many protesters held up flags in solidarity with African countries including Congo, Rwanda and Namibia.
Protesters also read out the names of scores of people killed in Gaza over the last 14 months.
The group called for an immediate ceasefire, an end to the illegal occupation, and the recognition of the Palestinian people ' s rights.
This protest comes days after Amnesty International found"sufficient basis"to conclude that"Israel committed and is continuing to commit genocide against Palestinians in the occupied Gaza Strip".
In a statement, the Palestine Solidarity Campaign ( PSC ), said :"We stand in solidarity with Palestinians and all other victims of genocide, both past and present, including Indigenous Americans, Indigenous Australians, Congolese, Sudanese, Rwandans, Namibians, Irish, Bosnians, Poles, Jews, Roma, Czechs, Armenians, and Slavs."We hold Israel, with its US backing, accountable for its ongoing genocide against the Palestinian people in Gaza.
Israel ' s continued forced expulsions in the West Bank and Gaza, along with its occupation and violence, represent a clear violation of international human rights, including the prohibition of genocide."Usuf Chikte, PSC Coordinator, told GroundUp :"The genocide in Gaza is being live-streamed.
Nobody in this world can say that they didn't know."In South Africa, we have a special obligation both morally and legally to ensure that apartheid never happens again.""We are reading the names of those who have been killed because they won't be forgotten.
We want to say that we are all equal in life and in death,"he said.
Earlier in the day, health workers held a press conference in District Six, demanding urgent action against Israel.
This coincided with similar events by health workers in Canada, USA, UK, Ireland, Turkey, Malaysia, and Namibia.
They called for an end to the targeting of healthcare workers, ambulance staff, aid workers in the field and journalists.
In a statement, Daniel Bloch, executive director of the Cape South African Jewish Board of Deputies, said he rejected"with utter contempt"the"baseless allegations of genocide"."A war which was started by Hamas on 7 October 2023, when they invaded Israel, is taking place in the Middle East, not a genocide.
However, we continue to deplore the loss of all life, we call for the release of all hostages and an end to the war.""The International Court of Justice, the primary judicial body of the United Nations, has made no ruling with regards to accusations of genocide, therefore any claims made by other organisations to the contrary, are misplaced."Bloch said"lies and mistruths"spread by anti-Israel movements"only enflame the hatred and intolerance towards the Jewish community, our supporters and \textbf{friends} and could well have led to the disgusting act of terror which occurred this past Friday 6 December, at our community building."He said an explosive device had been \textbf{thrown} into the complex"with the purpose of causing damage and creating fear and intimidation"."This investigation was concluded Monday morning by the bomb disposal unit.
The case docket has now been assigned to the Directorate for Priority Investigation Hawks for further investigation.
Despite this attack, our community remains resilient, united and strong and we will not be intimidated by these haters."News24 quoted police spokesperson Colonel Andr?
Traut saying that the police had found a possible explosive device which had not detonated and had removed it from the scene.
He confirmed that the docket had been assigned to the Hawks.
According to Al Jazeera, more than 44,000 Palestinians have been killed since October last year.
At least 2,000 people have also been killed in Lebanon.
About 1,100 Israelis were killed in the Hamas attack on Israel on 7 October.
Protesters held up flags in solidarity with African countries including Congo, Rwanda and Namibia.
You may republish this article, so long as you credit the authors and GroundUp, and do not change the text.
Please include a link back to the original article.
We put an invisible pixel in the article so that we can count traffic to republishers.
All analytics tools are solely on our servers.
We do not give our logs to any third party.
Logs are deleted after two weeks.
We do not use any IP address identifying information except to count regional traffic.
We are solely interested in counting hits, not tracking users.
If you republish, please do not delete the invisible pixel.

% matched lemmas: commemoration, friend, paint, throw
\end{textsample}
